۴ ۵

او از روی زمین دستش را بلند کرد و به صندلی جلویی کوچک با روکش قرمز لکه‌دار مشتی کوبید. هر از گاهی، لرزشی از خشم روی صورتش یا جاری شدن اشک‌هایی از چشمانش را می‌دیدم. امیدوار بودم تا زمانی که بتوانم به او کمک کنم، آرام بماند.

صدای آژیر نوحه‌اش همچنان بلندتر می‌شد، و اتاق را پر از صدای مردی می‌کرد که تمامی خصوصیات متمایزکننده انسان از کودک و حیوان را از دست داده است. او با خونسردی گفت: «می‌خواهم هنرهای رزمی یاد بگیرم تا وقتی که بخواهم کسی را بکشم، بتوانم چیزی درباره‌اش بکنم.»

او یک روپوش ابریشمی طلایی به تن داشت که چند سایز برایش کوچک بود و زانوهای زخمی‌اش را به نمایش می‌گذاشت. انتهای کمربند روپوش به سختی به هم می‌رسید تا یک گره بسازد و پرده‌های روپوش نیم‌متری از هم باز بودند، و سینه‌ای رنگ‌پریده و بی‌مو و شورت‌ای کرم‌رنگ از برند Calvin Klein را نشان می‌دادند.

تنها چیز دیگری که روی بدن لرزانش بود، یک کلاه زمستانی بود که محکم روی جمجمه‌اش کشیده شده بود. ماه ژوئن در لس‌آنجلس بود.

او دوباره صحبت می‌کرد: «این زندگی پوچ است.» او با چشمان خیس و قرمز به من نگاه کرد و گفت: «مثل بازی دوز است. هیچ راهی برای بردن نیست، پس بهترین کار این است که اصلاً بازی نکنی.»

در خانه کس دیگری نبود. باید با این وضعیت کنار می‌آمدم. باید قبل از اینکه او از حالت گریه به خشم برگردد، او را آرام کنم. هر دور از احساسات بدتر می‌شد و این بار می‌ترسیدم کاری کند که جبران‌پذیر نباشد.

نمی‌توانستم اجازه دهم Mystery در حضور من بمیرد. او بیش از یک دوست بود؛ او یک مربی بود. او زندگی من و زندگی هزاران نفر دیگر مثل من را تغییر داده بود. باید به او والیوم، زاناکس، ویکودین، هر چیزی که می‌توانستم، می‌دادم.

دفترچه تلفنم را برداشتم و صفحاتی را برای یافتن کسانی که احتمالاً قرص دارند، می‌گذراندم — افرادی مانند اعضای گروه‌های موسیقی راک، زنانی که تازه جراحی پلاستیک انجام داده‌اند یا بازیگران کودک سابق. اما هر کسی که تماس گرفتم در خانه نبود، دارویی نداشت یا مدعی نبودن دارو بود چون نمی‌خواست به اشتراک بگذارد.

فقط یک نفر باقی مانده بود: زنی که باعث سقوط Mystery شده بود. او یک دختر مهمانی بود؛ حتماً باید چیزی داشته باشد. کاتیا، روسی خردسال با صدای ظریف و انرژی یک توله سگ پامرانیان، در ده دقیقه جلو درب بود با یک قرص زاناکس و نگاهی نگران.

به او هشدار دادم: «داخل نیا، ممکن است او تو را بکشد.» نه اینکه او کاملاً لایقش نبود، البته. یا حداقل در آن زمان اینطور فکر می‌کردم. به Mystery قرص دادم و یک لیوان آب و منتظر ماندم تا گریه‌ها به خرخر تبدیل شوند. سپس او را با یک جفت چکمه سیاه، شلوار جین و تی‌شرت خاکستری کمک کردم. او حالا آرام بود، مانند یک بچه بزرگ.

به او گفتم: «می‌برمت جایی که کمکت کنند.» او را به بیرون بردم و توی کوروت زنگار گرفته‌ام جا دادم.

مقصد ما مرکز سلامت روانی هالیوود در خیابان واین بود. بلوکی زشت از بتن احاطه‌شده روز و شب توسط مردان بی‌خانمان که به تیرهای چراغ جیغ می‌زدند، زن‌پوش‌هایی که درون چرخ‌دستی‌های خرید زندگی می‌کردند، و بقیه انسان‌های مانده‌ای که در جایی که خدمات اجتماعی رایگان یافت می‌شد، اردو می‌زدند.

Mystery، فهمیدم یکی از آنها بود. او فقط کاریزما و استعدادی داشت که دیگران را به سوی خود می‌کشید و از تنها ماندن او در جهان جلوگیری می‌کرد.

او دارای دو ویژگی بود که تقریباً در هر ستاره راکی که تا به حال مصاحبه کرده بودم، دیده بودم: چشمانی با نگاه ک crazy و ناتوانی مطلق در اینکه برای خودش کاری انجام بدهد. او را به لابی آوردم، او را ثبت‌نام کردم و با هم منتظر نوبت خود با یکی از مشاوران شدیم.

او روی صندلی پلاستیکی ارزان‌قیمت سیاه نشسته بود و با حالتی ماتم زده به دیوارهای آبی موسسه خیره شد. یک ساعت گذشت. او شروع به بی‌قراری کرد.

دو ساعت گذشت. پیشانیش چین و چروک افتاد؛ چهره‌اش تاریک شد. سه ساعت گذشت. گریه‌ها آغاز شد.

چهار ساعت گذشت. او از صندلی‌اش پرید و از اتاق انتظار و از در ساختمان خارج شد. او به سرعت قدم برمی‌داشت، مثل مردی که می‌داند کجا می‌رود، اگرچه پروژه هالیوود سه مایل دورتر بود. او را در آن‌سوی خیابان دنبال کردم و در بیرون یک مرکز خرید کوچک به او رسیدم. بازویش را گرفتم و او را به آرامش به اتاق انتظار برگرداندم.

پنج دقیقه. ده دقیقه. بیست دقیقه. سی دقیقه. او دوباره بلند شد و رفت. او را دنبال کردم. دو مددکار اجتماعی بی‌فایده در لابی ایستاده بودند.

فریاد زدم: «او را متوقف کنید!» یکی از آنها گفت: «نمی‌توانیم، او ساختمان را ترک کرده است.»

پرسیدم: «پس شما فقط قصد دارید به یک مردی که گرایش به خودکشی دارد اجازه دهید اینجا را ترک کند؟» نمی‌توانستم زمان را با جروبحث هدر دهم. «فقط مطمئن شوید که یک درمانگر آماده دیدار با او باشد اگر توانستم او را برگردانم اینجا.»

از در بیرون رفتم و به سمت راستم نگاه کردم. او آنجا نبود. نگاه کردم